\section{Pruebas y CI/CD}
\label{sec:pruebas-cicd}

\subsection{Pirámide de pruebas}

La estrategia de pruebas del sistema cubre cuatro niveles. Las pruebas
unitarias emplean JUnit~5 y Mockito para el backend, aislando cada
componente de sus dependencias mediante dobles de prueba de los puertos
de dominio, sin necesitar infraestructura real. Las pruebas de integración verifican la persistencia contra instancias reales de
CockroachDB. En \texttt{auth-service} se emplea Testcontainers para levantar
dinámicamente una instancia aislada durante la ejecución de las pruebas,
mientras otros escenarios de integración utilizan la infraestructura definida
en Docker Compose. Las pruebas de contrato,
implementadas con Pact, validan la compatibilidad entre los
microservicios consumidores y proveedores de la API interna. Las
pruebas de extremo a extremo emplean Playwright para la aplicación web,
en tres motores de navegador (Chromium, Firefox y WebKit), y Espresso
para las pruebas instrumentadas de la aplicación móvil. Finalmente, las
pruebas de carga con Locust evalúan el comportamiento del sistema bajo
tráfico simulado.

\subsection{Cobertura de código por microservicio}

La cobertura de líneas se mide con JaCoCo en cada microservicio backend,
con un umbral mínimo exigido del 70\,\% verificado automáticamente
mediante la regla \texttt{jacoco:check} en el ciclo de vida de Maven.
La Tabla~\ref{tab:cobertura} resume el estado medido a la fecha de este
informe.

\begin{table}[htbp]
\centering
\caption{Cobertura de código por microservicio (JaCoCo)}
\label{tab:cobertura}
\begin{tabular}{lll}
\toprule
Microservicio & Cobertura & Tests \\
\midrule
academico-laboratorios-service & $\geq$ 70\,\% & 254 \\
usuarios-service & Por confirmar & --- \\
reservas-solicitudes-service & $>80\,\%$ líneas; $>48\,\%$ ramas & 110 \\
auth-service & $\geq$ 70\,\% & 30 \\
api-gateway & 78.70\,\% & 11 \\
\bottomrule
\end{tabular}
\end{table}

El microservicio \texttt{academico-laboratorios-service} alcanza una
cobertura verificada superior al umbral exigido, con 254 pruebas
unitarias distribuidas en las cuatro capas de su arquitectura
hexagonal: 11 servicios de aplicación, 12 controladores de
presentación, la fachada de coordinación (\texttt{LaboratorioDetalleFacadeImpl})
y los 11 adaptadores de persistencia que implementan el patrón
Repository sobre Spring Data JPA.

Para \texttt{reservas-solicitudes-service}, la verificación final del
módulo ejecutó 110 pruebas sin fallos, errores ni pruebas omitidas. La
suite cubre ownership del docente, ámbito del administrador de piso,
transiciones y propuestas, cancelación transaccional, idempotencia de
creación y aprobación, migraciones Flyway y concurrencia sobre
CockroachDB. JaCoCo registró más del 80\,\% de líneas y más del
48\,\% de ramas; el umbral se comprueba durante \texttt{mvn verify} para
detectar regresiones futuras.

En \texttt{auth-service}, la suite completa ejecutó 30 pruebas sin fallos,
errores ni omisiones. Además, se incorporó una prueba de integración del
repositorio utilizando Testcontainers con una instancia real de CockroachDB,
sobre la cual Flyway aplica las migraciones antes de validar operaciones de
persistencia y lectura mediante JPA.

En \texttt{api-gateway}, JaCoCo registró una cobertura de líneas del
78.70\,\% (85 de 108 líneas), superando el umbral mínimo del 70\,\%.
La ejecución de \texttt{mvn clean verify} completó 11 pruebas sin fallos
ni errores, y la regla \texttt{jacoco:check} quedó integrada en la fase
\texttt{verify} para impedir regresiones de cobertura.

En la web, las pruebas unitarias, ESLint y la construcción de producción
finalizaron correctamente. El flujo E2E de reservas obtuvo dos ejecuciones
consecutivas de 2/2 casos aprobados, sin reintentos automáticos, y la suite
E2E preexistente completó 6/6 casos. En Android se ejecutaron 76 pruebas
sin fallos ni omisiones, Android Lint terminó con cero errores y se generó
el APK de depuración. La validación visual en dispositivo quedó fuera de
la evidencia automatizada por no existir emulador o dispositivo disponible
en el entorno de ejecución.

Las pruebas de contrato con Pact incluyen además el flujo de autenticación
\texttt{POST /api/v1/auth/login}. El contrato valida la estructura de la
solicitud, la respuesta HTTP exitosa, los tokens de acceso y renovación,
su tiempo de expiración y la información del usuario autenticado, incluidos
roles, permisos y tipos de perfil. La suite de contratos ejecutó seis pruebas
sin fallos ni errores.

\subsection{Pipeline de integración continua}

El pipeline de \texttt{GitHub Actions}, definido en
\texttt{.github/workflows/ci.yml}, ejecuta un conjunto de trabajos
encadenados sobre cada \emph{push} y \emph{pull request} a la rama de
trabajo, siguiendo el enfoque de entrega continua descrito por
Humble y Farley \cite{humble2010continuous}.

\begin{table}[htbp]
\centering
\caption{Jobs del pipeline de integración continua}
\label{tab:jobs-cicd}
\begin{tabular}{lll}
\toprule
Job & Propósito & Estado \\
\midrule
lint & Checkstyle (backend) + ESLint (web) & Completo \\
build-test & Compilación y pruebas unitarias Maven & Completo \\
crdb-tests & Pruebas contra el clúster CockroachDB E3 & Pendiente \\
test-web & Pruebas unitarias con Vitest & Completo \\
test-mobile & Pruebas unitarias e instrumentadas & Completo \\
build-mobile-apk & Empaquetado del APK de depuración & Completo \\
build-images & Publicación de imágenes en GHCR & Completo \\
integration & Levantamiento y verificación del stack & Completo \\
\bottomrule
\end{tabular}
\end{table}

El job \texttt{crdb-tests} permanece como una tarea pendiente
explícitamente marcada en el propio flujo de trabajo, condicionada a la
finalización del esquema de \texttt{reservas-solicitudes-service}; se
documenta con honestidad en lugar de presentarlo como resuelto.

\subsection{Análisis estático de código}

El job \texttt{lint} integra dos herramientas de análisis estático
según la tecnología de cada componente. Para los microservicios backend
en Java se emplea \texttt{maven-checkstyle-plugin} con una configuración
propia (\texttt{checkstyle.xml}) que cubre reglas de importaciones no
utilizadas, llaves obligatorias, longitud máxima de línea (120
caracteres) y equivalencia entre \texttt{equals} y \texttt{hashCode}.
Para la aplicación web se emplea \texttt{ESLint} sobre el código
TypeScript/React. Ambas herramientas se ejecutan automáticamente en
cada integración y detienen el pipeline si detectan violaciones.
